% Rules of the underlying linear and affine algebras (GLA, GAA) that are used
% as background throughout, but are not contributions of this thesis.
\medskip
\noindent
\fbox{\begin{minipage}{0.97\linewidth}
    \begin{center}\textbf{(G1) Matrices}\quad{\footnotesize\cite{bonchi2017interacting,paixao2022highlevel}}\end{center}
    \begin{center}\begin{tabular}{r c c c}
      \diagrow{(comp.)}{
          \pic[val=A] (g1) at (0.500,0.500) {matrix};
          \pic[val=B] (g2) at (1.750,0.500) {matrix};
          \draw (g1-out-0) to[out=0,in=180] (g2-in-0);
      }{
          \pic[val=BA] (g1) at (0.500,0.500) {matrix};
      }
      \diagrow{(row)}{
          \pic[val=A] (g1) at (0.500,1.000) {matrix};
          \pic[val=B] (g2) at (0.500,0.000) {matrix};
          \pic (g3) at (1.750,0.500) {multiplication};
          \draw (g1-out-0) to[out=0,in=180] (g3-in-0);
          \draw (g2-out-0) to[out=0,in=180] (g3-in-1);
      }{
          \pic[val={[A \ B]}] (g1) at (1.000,0.000) {2dmatrix};
          \draw (0.000,0.300) -- (g1-in-0);
          \draw (0.000,-0.300) -- (g1-in-1);
      }
      \diagrow{(antipode)}{
          \pic (g1) at (0.500,0.500) {antipode};
      }{
          \pic[val=-1] (g1) at (0.500,0.500) {scalar};
      }
      \diagrow{(zero sc.)}{
          \pic[val=0] (g1) at (0.500,0.500) {scalar};
      }{
          \pic (g1) at (0.500,0.500) {discard};
          \pic (g2) at (1.750,0.500) {zero};
          \draw (0.000,0.500) -- (g1-in-0);
          \draw (g2-out-0) -- (2.500,0.500);
      }
    \end{tabular}\end{center}
\end{minipage}}

\medskip
\noindent
\fbox{\begin{minipage}{0.97\linewidth}
    \begin{center}\textbf{(G2) Transposition and cospan form}\quad{\footnotesize\cite{bonchi2017interacting}}\end{center}
    A diagram may be slid along a cap, exchanging it for its transpose; on
    $\mathbf{GLA}$ diagrams the transpose is the converse relation, so a matrix
    becomes a comatrix.
    \begin{center}\begin{tabular}{r c c c}
      \diagrow{(slide)}{
          \pic[val=A] (g1) at (0.500,1.000) {matrix};
          \draw (g1-out-0) -- (1.600,1.000) arc (90:-90:0.500) -- (0.000,0.000);
      }{
          \pic[val=A] (g1) at (0.500,0.000) {comatrix};
          \draw (g1-out-0) -- (1.600,0.000) arc (-90:90:0.500) -- (0.000,1.000);
      }
      \diagrow{(cospan)}{
          \pic[val=g] (g1) at (0.500,0.500) {relation};
      }{
          \pic[val=E] (g1) at (0.500,0.500) {matrix};
          \pic[val=F] (g2) at (1.750,0.500) {comatrix};
          \draw (g1-out-0) to[out=0,in=180] (g2-in-0);
      }
    \end{tabular}\end{center}
\end{minipage}}

\medskip
\noindent
\fbox{\begin{minipage}{0.97\linewidth}
    \begin{center}\textbf{(G3) Bending}\quad{\footnotesize\cite{bonchi2017interacting,bonchi2019affine}}\end{center}
    The white cap is the black cap with an antipode on one leg; this is what
    turns an equation into a pair of bendable wires.
    \begin{center}\begin{tabular}{r c c c}
      \diagrow{(white cap)}{
          \pic (g1) at (0.500,0.500) {multiplication};
          \pic (g2) at (1.750,0.500) {cozero};
          \draw (g1-out-0) to[out=0,in=180] (g2-in-0);
      }{
          \pic (g1) at (0.500,0.000) {antipode};
          \draw (g1-out-0) -- (1.600,0.000) arc (-90:90:0.500) -- (0.000,1.000);
      }
      \diagrow{(unbend $\oplus$)}{
          \pic (g1) at (0.500,0.500) {comultiplication};
          \pic (g2) at (1.750,0.000) {antipode};
          \draw (g1-out-0) -- (2.500,1.000);
          \draw (g1-out-1) to[out=0,in=180] (g2-in-0);
          \draw (g2-out-0) -- (2.500,0.000);
      }{
          \pic (g1) at (1.000,0.500) {multiplication};
          \draw (0.000,1.000) to[out=0,in=180] (g1-in-0);
          \draw (g1-out-0) -- (1.900,0.500) arc (90:-90:0.500) -- (0.000,-0.500);
      }
    \end{tabular}\end{center}
\end{minipage}}

\medskip
\noindent
\fbox{\begin{minipage}{0.97\linewidth}
    \begin{center}\textbf{(G4) The black structure is a spider}\quad{\footnotesize\cite{bonchi2017interacting}}\end{center}
    Any connected network of black nodes with $a$ inputs and $b$ outputs equals
    the single node $\mathsf{cc}_a;\mathsf{cp}_b$; the two laws below are the
    cases used here.
    \begin{center}\begin{tabular}{r c c c}
      \diagrow{(special)}{
          \pic (g1) at (0.500,0.500) {copy};
          \pic (g2) at (1.750,0.500) {cocopy};
          \draw (g1-out-0) to[out=0,in=180] (g2-in-0);
          \draw (g1-out-1) to[out=0,in=180] (g2-in-1);
      }{
          \draw (0.000,0.500) -- (1.000,0.500);
      }
      \diagrow{(merge/del)}{
          \pic (g1) at (0.500,0.500) {cocopy};
          \pic (g2) at (1.750,0.500) {discard};
          \draw (g1-out-0) -- (g2-in-0);
      }{
          \pic (g1) at (0.500,1.000) {discard};
          \pic (g2) at (0.500,0.000) {discard};
          \draw (0.000,1.000) -- (g1-in-0);
          \draw (0.000,0.000) -- (g2-in-0);
      }
      \diagrow{(split/un)}{
          \pic (g1) at (0.500,0.500) {codiscard};
          \pic (g2) at (1.750,0.500) {copy};
          \draw (g1-out-0) -- (g2-in-0);
      }{
          \pic (g1) at (0.500,1.000) {codiscard};
          \pic (g2) at (0.500,0.000) {codiscard};
          \draw (g1-out-0) -- (1.250,1.000);
          \draw (g2-out-0) -- (1.250,0.000);
      }
    \end{tabular}\end{center}
\end{minipage}}
